% Relativistic Chapter II, Sections 5--9
% Bénard problem: perturbation equations in a relativistic BDNK framework
% Based on Chandrasekhar Ch.~II §5--9, extended to special-relativistic hydrodynamics
% Conventions: see BDNK_CONVENTIONS.md

\chapter{RELATIVISTIC THERMAL INSTABILITY OF A HEATED FLUID LAYER}
\label{ch:rel2}

\begin{center}
  {\Large 1. THE RELATIVISTIC B\'ENARD PROBLEM}
\end{center}

%======================================================================
\section{Introduction}\label{sec:rel5}
%======================================================================

The classical theory of thermal instability in a horizontal layer of fluid
heated from below, as developed by Rayleigh and expounded in Chapter~II,
rests on the Navier--Stokes equations supplemented by Fourier's law of
heat conduction.  While this framework is entirely adequate for ordinary
laboratory fluids at non-relativistic temperatures, there are compelling
reasons---both of principle and of astrophysical application---to formulate
the B\'enard problem within a relativistic setting.

\emph{First}, Fourier's law predicts infinite speed of propagation for
thermal signals: a temperature perturbation at one point is felt
instantaneously everywhere.  This parabolic character of the
diffusion equation violates relativistic causality.  The BDNK
(Bemfica--Disconzi--Noronha--Kovtun) formalism of first-order causal
viscous relativistic hydrodynamics addresses this issue not by introducing
relaxation equations (as in the Israel--Stewart approach), but by a careful
choice of \emph{hydrodynamic frame} and constrained first-order transport
coefficients.  The resulting theory is causal, strongly hyperbolic, and
involves \emph{no additional dynamical degrees of freedom} beyond the
standard hydrodynamic variables $(\varepsilon, n, u^\mu)$.

\emph{Second}, in relativistic fluids it is not the rest-mass density
$\rho_0$ alone that provides the inertia; rather, it is the \emph{enthalpy
density} $w = (\varepsilon + p)/c^2$ (where $\varepsilon$ is the energy density
and $p$ the pressure) that plays the role of inertial mass density.
This distinction becomes significant in astrophysical plasmas---neutron-star
oceans, accretion-disc atmospheres, the early-universe quark--gluon
plasma---where thermal and binding energies are a non-negligible fraction
of the rest-mass energy.

\emph{Third}, a systematic relativistic treatment reveals how the
classical Rayleigh number is modified when pressure contributions to
inertia, relativistic thermodynamics, and causal heat transport are
included.  The corrections are of order $v^2/c^2$, and they qualitatively
change the stability boundary in the high-temperature regime.

Our aim in the present chapter is therefore to derive, from first principles
of relativistic fluid dynamics in the BDNK formalism, the perturbation
equations governing the onset of convective instability in a relativistic
fluid layer heated from below.  The analysis parallels that of
Chandrasekhar's classical treatment (\S\S~5--9 of Chapter~II) step by step,
so that the non-relativistic limit can be recovered at every stage by
taking $c \to \infty$.

%======================================================================
\section{The nature of the physical problem}\label{sec:rel6}
%======================================================================

Consider a horizontal layer of relativistic fluid of depth $d$, confined
between two plane boundaries at $z = 0$ (bottom) and $z = d$ (top).
The bottom boundary is maintained at temperature $T_0$ and the top at
$T_1 < T_0$, so that a steady adverse temperature gradient
$\beta = (T_0 - T_1)/d > 0$ is established.  Gravity acts in the
$-z$ direction with (proper) acceleration~$g$.

In the classical theory, the inertia of a fluid element is determined
solely by the rest-mass density $\rho$.  In the relativistic theory, the
\emph{total inertial mass density} is
\begin{equation}
\label{eq:rel2-1}
  w \;=\; \frac{\varepsilon + p}{c^2}\,,
\end{equation}
where $\varepsilon$ is the total energy density (rest-mass energy plus
internal energy) and $p$ is the isotropic pressure.  For a
non-relativistic ideal gas, $\varepsilon \approx \rho_0 c^2$ and
$p \ll \rho_0 c^2$, so $w \approx \rho_0$; but in general the
pressure and thermal energy contribute to $w$ and thereby modify the
buoyancy force.

The thermal expansion of the fluid is described, as classically, by
\begin{equation}
\label{eq:rel2-2}
  \rho = \rho_0\bigl[1 - \alpha(T - T_0)\bigr],
\end{equation}
where $\alpha$ is the coefficient of volume expansion.  However, the energy
density $\varepsilon$ and the pressure $p$ each depend on temperature, so
the enthalpy density $w$ varies as
\begin{equation}
\label{eq:rel2-3}
  w = w_0\bigl[1 - \alpha_w(T - T_0)\bigr], \qquad
  \alpha_w = -\frac{1}{w_0}\frac{\partial w}{\partial T}\bigg|_{T_0}\,,
\end{equation}
where $\alpha_w$ is an \emph{effective enthalpy expansion coefficient}.  In
general $\alpha_w \neq \alpha$; the discrepancy encodes genuinely
relativistic corrections to the buoyancy.

The instability is still controlled by a competition between buoyancy
(proportional to $g\,\alpha\,\beta\,d$) and dissipation (viscous and thermal).
Following the classical pattern, we define the \emph{relativistic Rayleigh number}
\begin{equation}
\label{eq:rel2-4}
  \mathrm{Ra}_{\mathrm{rel}} = \frac{g\,\alpha\,\beta\,d^4}{\kappa_T\,\nu}
  \left[1 + \mathcal{O}\!\left(\frac{v^2}{c^2}\right)\right],
\end{equation}
where $\kappa_T$ is the relativistic thermal diffusivity (to be defined below)
and $\nu = \eta_{\mathrm{s}} / w_0$ is the \emph{relativistic} kinematic viscosity
(note: the inertial mass density $w_0$, not $\rho_0$, enters the denominator).
Instability sets in when $\mathrm{Ra}_{\mathrm{rel}}$ exceeds a critical value
$\mathrm{Ra}_{\mathrm{rel},c}$.

A key qualitative difference from both the classical problem and the
Israel--Stewart formulation is the mechanism by which causality is enforced.
In the BDNK framework, there are \emph{no relaxation times}
$\tau_q$, $\tau_\pi$; instead, causality and hyperbolicity are ensured by
the choice of hydrodynamic frame, encoded in frame coefficients that
modify the out-of-equilibrium definitions of energy density, pressure,
and heat flux.  The perturbation equations are therefore \emph{lower order}
in time derivatives than in the Israel--Stewart theory, yet remain
hyperbolic.

%======================================================================
\section{The basic relativistic hydrodynamic equations}\label{sec:rel7}
%======================================================================

We work in the framework of special relativity in the presence of a
uniform gravitational field $g$ (i.e.\ the weak-field limit where the
spacetime is effectively Minkowski with a static gravitational potential
$\Phi = g\,z$).  Greek indices $\mu,\nu,\ldots = 0,1,2,3$ label spacetime
components; Latin indices $i,j,\ldots = 1,2,3$ label spatial components.
The metric signature is $(-,+,+,+)$, and the 4-velocity $u^\mu$ is
normalized according to
\begin{equation}
\label{eq:rel2-5}
  u^\mu u_\mu = -c^2.
\end{equation}
The projection tensor orthogonal to $u^\mu$ is
\begin{equation}
\label{eq:rel2-6}
  \Delta^{\mu\nu} = g^{\mu\nu} + \frac{u^\mu u^\nu}{c^2}\,.
\end{equation}

%----------------------------------------------------------------------
\subsection{The energy--momentum tensor}\label{sec:rel7a}
%----------------------------------------------------------------------

In the BDNK formalism, the energy--momentum tensor is decomposed in a
\emph{general hydrodynamic frame} (neither Landau nor Eckart):
\begin{equation}
\label{eq:rel2-7}
  T^{\mu\nu} = E\,\frac{u^\mu u^\nu}{c^2}
    + P\,\Delta^{\mu\nu}
    + \Pi^{\mu\nu}
    + \frac{1}{c^2}\bigl(Q^\mu u^\nu + Q^\nu u^\mu\bigr),
\end{equation}
where $E$ is the (out-of-equilibrium) energy density, $P$ is the
(out-of-equilibrium) isotropic pressure, $\Pi^{\mu\nu}$ is the
anisotropic stress tensor (symmetric, traceless, and orthogonal to
$u^\mu$), and $Q^\mu$ is the heat-flux 4-vector (orthogonal to $u^\mu$:
$Q^\mu u_\mu = 0$).

The crucial distinction from the Israel--Stewart decomposition is that
$E$, $P$, $Q^\mu$, and $\Pi^{\mu\nu}$ are given by \emph{algebraic}
(first-order gradient) constitutive relations, not by separate evolution
equations.  Specifically, the BDNK constitutive relations are:
\begin{align}
\label{eq:rel2-7a}
  E &= \varepsilon + \varepsilon_1, \qquad
  \varepsilon_1 = -\zeta_1\,\nabla_\alpha u^\alpha
    - \frac{\beta_1}{T}\,u^\alpha\partial_\alpha T
    - \alpha_1\,u^\alpha\partial_\alpha\!\left(\frac{\mu}{T}\right), \\[4pt]
\label{eq:rel2-7b}
  P &= p + \Pi_{\mathrm{bulk}}, \qquad
  \Pi_{\mathrm{bulk}} = -\zeta\,\nabla_\alpha u^\alpha + \ldots
    \;\text{(frame corrections)}, \\[4pt]
\label{eq:rel2-7c}
  Q^\mu &= -\kappa\,T\,\Delta^{\mu\alpha}\partial_\alpha\!\left(\frac{\mu}{T}\right)
    + \text{(BDNK frame corrections)}, \\[4pt]
\label{eq:rel2-7d}
  \Pi^{\mu\nu} &= -2\eta\,\sigma^{\mu\nu},
\end{align}
where $\varepsilon$ and $p$ are the \emph{equilibrium} energy density and
pressure, $\eta$ is the shear viscosity, $\zeta$ is the bulk viscosity,
$\kappa$ is the thermal conductivity, and $\varepsilon_1$, $\beta_1$,
$\alpha_1$, $\zeta_1$ are BDNK \emph{frame coefficients} that encode the
departure of the hydrodynamic frame from the Landau frame.  The shear
tensor $\sigma^{\mu\nu}$ is defined in equation~\eqref{eq:rel2-16} below.

\medskip
\noindent\textbf{Remark.}\quad
In the Landau frame one sets $\varepsilon_1 = 0$ and $Q^\mu = 0$, but
this choice renders the first-order theory acausal (the Eckart/Hiscock--Lindblom
instability).  The BDNK frame avoids this pathology by retaining
non-zero frame coefficients, which contribute to the principal part of the
equations of motion and restore strong hyperbolicity.

%----------------------------------------------------------------------
\subsection{The equation of continuity (baryon-number conservation)}%
\label{sec:rel7b}
%----------------------------------------------------------------------

The conservation of baryon number is expressed as
\begin{equation}
\label{eq:rel2-8}
  \nabla_\mu\!\bigl(\rho_0\, u^\mu\bigr) = 0,
\end{equation}
where $\rho_0$ is the rest-mass density (baryon mass density) measured in the
local rest frame.  In component form, in a coordinate system where the
Christoffel symbols vanish (Minkowski background), this reads
\begin{equation}
\label{eq:rel2-9}
  \frac{\partial}{\partial t}\bigl(\gamma\rho_0\bigr)
  + \frac{\partial}{\partial x^j}\bigl(\gamma\rho_0\, v^j\bigr) = 0,
\end{equation}
where $v^j$ is the 3-velocity, $\gamma = (1 - v^2/c^2)^{-1/2}$ is the
Lorentz factor, and $u^\mu = \gamma(c, v^j)$.

In the non-relativistic limit ($\gamma \to 1$, $v \ll c$), equation~\eqref{eq:rel2-9}
reduces to the classical continuity equation~(2) of Chandrasekhar's
Chapter~II.

%----------------------------------------------------------------------
\subsection{The equations of motion (momentum conservation)}%
\label{sec:rel7c}
%----------------------------------------------------------------------

The conservation of energy and momentum is encoded in
\begin{equation}
\label{eq:rel2-10}
  \nabla_\nu T^{\mu\nu} = f^\mu_{\mathrm{ext}},
\end{equation}
where $f^\mu_{\mathrm{ext}}$ is the external 4-force density.  For a uniform
gravitational field $g$ in the $-z$ direction,
\begin{equation}
\label{eq:rel2-11}
  f^\mu_{\mathrm{ext}} = -w\,g\,\lambda^\mu, \qquad
  \lambda^\mu = (0,0,0,1),
\end{equation}
where $\lambda^\mu$ is the unit vector in the $z$-direction (with appropriate
covariant components in the weak-field limit).

In the BDNK formalism, no separate evolution equations for $Q^\mu$ or
$\Pi^{\mu\nu}$ are needed: the dissipative fluxes are expressed directly
in terms of gradients of the primary hydrodynamic variables through
the constitutive relations~\eqref{eq:rel2-7a}--\eqref{eq:rel2-7d}.
Projecting equation~\eqref{eq:rel2-10} orthogonally to $u^\mu$ (using
$\Delta^{\mu}{}_{\alpha}$) yields the relativistic Euler equation with
first-order dissipative corrections:
\begin{equation}
\label{eq:rel2-12}
  w\, u^\nu\nabla_\nu u^\mu
  = -\Delta^{\mu\nu}\nabla_\nu p
    - \Delta^{\mu\nu}\nabla_\nu \Pi_{\mathrm{bulk}}
    + \Delta^{\mu}{}_{\alpha}\nabla_\beta\Pi^{\alpha\beta}
    + \frac{1}{c^2}\Delta^{\mu}{}_{\alpha}\bigl(
        Q^\alpha\nabla_\beta u^\beta
      + Q^\beta\nabla_\beta u^\alpha
      \bigr)
    + f^\mu_{\perp},
\end{equation}
where $w = (\varepsilon + p)/c^2$ is the enthalpy per unit volume divided by
$c^2$ (the relativistic inertial mass density), and
$f^\mu_\perp = \Delta^{\mu}{}_\alpha f^\alpha_{\mathrm{ext}}$
is the spatial projection of the gravitational force.

Since $\Pi^{\mu\nu} = -2\eta\,\sigma^{\mu\nu}$ algebraically in the BDNK
theory (no relaxation equation), the shear-stress divergence can be
directly expressed in terms of velocity gradients.

In the non-relativistic limit, $w \to \rho_0$,
$\Delta^{\mu\nu}\nabla_\nu p \to \partial p / \partial x^i$, the
$Q$-dependent terms vanish, and the shear-stress divergence reduces to
$\eta\nabla^2 u_i$.  We thus recover the Navier--Stokes equation~(18) of
Chapter~II.

%----------------------------------------------------------------------
\subsection{The energy equation}\label{sec:rel7d}
%----------------------------------------------------------------------

Projecting equation~\eqref{eq:rel2-10} along $u_\mu$ gives the energy
equation in the local rest frame:
\begin{equation}
\label{eq:rel2-13}
  u^\mu\nabla_\mu\varepsilon
  + (\varepsilon + p)\,\nabla_\mu u^\mu
  = -\Pi_{\mathrm{bulk}}\,\nabla_\mu u^\mu
    - \Pi^{\mu\nu}\nabla_\mu u_\nu
    - \nabla_\mu Q^\mu
    - \frac{1}{c^2}\,Q^\mu u^\nu\nabla_\nu u_\mu\,.
\end{equation}
The left-hand side is the relativistic generalization of the material
derivative of the internal energy plus the $p\,\mathrm{div}\,\mathbf{u}$
work term.  The right-hand side collects the dissipative contributions:
bulk-viscous heating, shear-viscous heating, heat-flux divergence, and
a relativistic correction involving the acceleration of the heat flux.

In the BDNK theory, the heat flux $Q^\mu$ appearing in
equation~\eqref{eq:rel2-13} is given by the constitutive
relation~\eqref{eq:rel2-7c}; no separate evolution equation for $Q^\mu$ is
needed.  The divergence $\nabla_\mu Q^\mu$ thus yields second-order spatial
derivatives of the temperature, producing the standard diffusive operator
at leading order.

In the non-relativistic limit, equation~\eqref{eq:rel2-13} reduces to
the heat equation~(39) of Chapter~II (upon identifying
$\varepsilon \approx \rho_0 c_v T + \mathrm{const}$ and
$\nabla_\mu Q^\mu \to -\partial/\partial x_j(k\,\partial T/\partial x_j)$).

%----------------------------------------------------------------------
\subsection{BDNK constitutive relations (causal first-order dissipation)}%
\label{sec:rel7e}
%----------------------------------------------------------------------

The crucial feature of the BDNK formalism is that it replaces the
Israel--Stewart \emph{relaxation equations} with \emph{algebraic}
constitutive relations at first order in gradients, while maintaining
causality through the choice of hydrodynamic frame.

\paragraph{Heat flux.}
In the BDNK general frame, the heat flux is given by
\begin{equation}
\label{eq:rel2-14}
  Q^\mu = -\kappa\,T\,\Delta^{\mu\alpha}\partial_\alpha\!\left(\frac{\mu}{T}\right)
    + \varepsilon_Q\,\Delta^{\mu\alpha}\,u^\beta\nabla_\beta u_\alpha
    + \text{(additional frame terms)},
\end{equation}
where $\kappa$ is the thermal conductivity, $\mu$ is the chemical potential,
and $\varepsilon_Q$ is a BDNK frame coefficient.  For a single-component
fluid, $\partial_\alpha(\mu/T)$ can be re-expressed in terms of $\partial_\alpha T$
and $\partial_\alpha p$ using the Gibbs--Duhem relation; in the simplest
case (no conserved charge gradients) this reduces to
\begin{equation}
\label{eq:rel2-14b}
  Q^\mu \approx -\kappa_{\mathrm{eff}}\,\Delta^{\mu\alpha}\partial_\alpha T
    + \text{(frame corrections)},
\end{equation}
where $\kappa_{\mathrm{eff}}$ absorbs thermodynamic factors.
Equation~\eqref{eq:rel2-14} is \emph{algebraic}---it contains no time
derivatives of $Q^\mu$ itself.  There is no relaxation time $\tau_q$.

\paragraph{Shear stress.}
The anisotropic stress is given by
\begin{equation}
\label{eq:rel2-15}
  \Pi^{\mu\nu} = -2\eta\,\sigma^{\mu\nu},
\end{equation}
where $\eta$ is the shear viscosity and
\begin{equation}
\label{eq:rel2-16}
  \sigma^{\mu\nu} = \frac{1}{2}\bigl(
    \Delta^{\mu\alpha}\nabla_\alpha u^\nu
    + \Delta^{\nu\alpha}\nabla_\alpha u^\mu
  \bigr)
  - \frac{1}{3}\Delta^{\mu\nu}\nabla_\alpha u^\alpha
\end{equation}
is the shear tensor.  This is again an algebraic relation; there is no
relaxation time $\tau_\pi$.  In the non-relativistic Boussinesq limit,
$\Pi_{ij} = -\eta(\partial_i u_j + \partial_j u_i)$, recovering the
Navier--Stokes stress.

\paragraph{Bulk viscous pressure.}
The bulk viscous pressure is
\begin{equation}
\label{eq:rel2-17}
  \Pi_{\mathrm{bulk}} = -\zeta\,\nabla_\mu u^\mu,
\end{equation}
where $\zeta$ is the bulk viscosity.  Again, no relaxation time
$\tau_\Pi$ appears.

\paragraph{Comparison with Israel--Stewart.}
In the Israel--Stewart theory, $q^\mu$, $\pi^{\mu\nu}$, and $\Pi$ are
promoted to independent dynamical variables satisfying relaxation equations
(e.g.\ $\tau_q\,u^\alpha\nabla_\alpha q^\mu + q^\mu = \ldots$).  The BDNK
approach eliminates these extra degrees of freedom entirely: the
conservation laws $\nabla_\mu T^{\mu\nu} = 0$ and $\nabla_\mu(\rho_0 u^\mu) = 0$,
together with the algebraic constitutive
relations~\eqref{eq:rel2-14}--\eqref{eq:rel2-17}, form a \emph{closed}
system for $(\varepsilon, u^\mu, \rho_0)$.  Causality is ensured not by
relaxation times but by the frame coefficients.

%----------------------------------------------------------------------
\subsection{The rate of dissipation}\label{sec:rel7f}
%----------------------------------------------------------------------

The relativistic analogue of the viscous dissipation function $\Phi$ of
equation~(29) in Chapter~II is
\begin{equation}
\label{eq:rel2-18}
  \Phi_{\mathrm{rel}} = \Pi^{\mu\nu}\sigma_{\mu\nu}
    + \Pi_{\mathrm{bulk}}\,\nabla_\mu u^\mu\,.
\end{equation}
Since $\Pi^{\mu\nu} = -2\eta\,\sigma^{\mu\nu}$ in the BDNK theory,
\begin{equation}
\label{eq:rel2-18b}
  \Phi_{\mathrm{rel}} = 2\eta\,\sigma^{\mu\nu}\sigma_{\mu\nu}
    + \zeta\,(\nabla_\mu u^\mu)^2\,.
\end{equation}
This quantity is non-negative by the second law of thermodynamics
($\Phi_{\mathrm{rel}} \ge 0$), provided the transport coefficients
$\eta \ge 0$, $\zeta \ge 0$.  In the non-relativistic limit
with an incompressible fluid, $\Phi_{\mathrm{rel}} \to 2\mu\,e_{ij}^2$,
recovering equation~(31) of Chapter~II.

%======================================================================
\section{The relativistic Boussinesq approximation}\label{sec:rel8}
%======================================================================

The classical Boussinesq approximation (\S~8 of Chapter~II) rests on the
smallness of the coefficient of thermal expansion $\alpha$: density
variations are neglected everywhere except in the gravitational body-force
term, where even a small $\delta\rho$ produces a significant buoyancy
acceleration.

In the relativistic setting, we adopt an analogous approximation but with
an important modification: the inertia of the fluid is governed by the
enthalpy density $w = (\varepsilon + p)/c^2$, not merely by $\rho_0$.
The \textbf{relativistic Boussinesq approximation} therefore consists of the
following prescriptions:

\begin{enumerate}
\item \textbf{Incompressibility:} The divergence of the 4-velocity in the
  local rest frame is negligible,
  \begin{equation}
  \label{eq:rel2-19}
    \nabla_\mu u^\mu \approx 0,
  \end{equation}
  which in the low-velocity limit gives $\partial u_j / \partial x_j = 0$
  as in the classical case (equation~(41) of Chapter~II).

\item \textbf{Constant inertial mass density:} In all terms except the
  gravitational force, we replace $w$ by its reference value~$w_0$:
  \begin{equation}
  \label{eq:rel2-20}
    w \approx w_0 \;=\; \frac{\varepsilon_0 + p_0}{c^2}\,.
  \end{equation}

\item \textbf{Buoyancy:} In the gravitational force term, we retain the
  \emph{full} variation of $w$:
  \begin{equation}
  \label{eq:rel2-21}
    \delta w \;=\; w - w_0
    \;=\; -w_0\,\alpha_w\,(T - T_0),
  \end{equation}
  where $\alpha_w$ is the enthalpy expansion coefficient defined in
  equation~\eqref{eq:rel2-3}.  Crucially, $\alpha_w$ receives contributions
  from both $\delta\varepsilon$ and $\delta p$:
  \begin{equation}
  \label{eq:rel2-22}
    \alpha_w = \alpha + \frac{1}{w_0 c^2}
    \left(\frac{\partial p}{\partial T}\right)_{\!\rho}\,,
  \end{equation}
  so that the pressure variation with temperature (negligible classically)
  contributes to the buoyancy in the relativistic regime.

\item \textbf{Constant transport coefficients:} The viscosity
  $\eta$, thermal conductivity $\kappa$, and BDNK frame
  coefficients are treated as constants evaluated at the reference state.
\end{enumerate}

Under these approximations, the relativistic momentum equation~%
\eqref{eq:rel2-12} becomes, in the low-velocity limit (3-vector form),
after substituting the BDNK algebraic constitutive relation
$\Pi_{ij} = -2\eta\,\sigma_{ij}$ directly:
\begin{equation}
\label{eq:rel2-23}
  w_0\!\left(\frac{\partial u_i}{\partial t}
    + u_j\frac{\partial u_i}{\partial x_j}\right)
  = -\frac{\partial\,\delta p}{\partial x_i}
    + g\,\alpha_w\,w_0\,\theta\,\lambda_i
    + \eta\,\nabla^2 u_i\,,
\end{equation}
where $\theta$ is the temperature perturbation.  Note that, unlike the
Israel--Stewart formulation, the shear stress has been substituted
\emph{algebraically}; no relaxation equation is needed.

The heat equation, under the Boussinesq approximation and using the
BDNK constitutive relation for the heat flux~\eqref{eq:rel2-14b}, yields
a \emph{first-order-in-time} diffusion equation for $\theta$ (see
\S~\ref{sec:rel9} for the linearized form), in contrast to the
second-order telegraph equation of the Israel--Stewart formulation.

Dividing equation~\eqref{eq:rel2-23} by $w_0$, we define the
\emph{relativistic kinematic viscosity}
\begin{equation}
\label{eq:rel2-24}
  \nu_{\mathrm{rel}} = \frac{\eta}{w_0}
  = \frac{\eta\,c^2}{\varepsilon_0 + p_0}\,.
\end{equation}
In the classical limit $\varepsilon_0 \to \rho_0 c^2$, $p_0 \ll \rho_0 c^2$,
so $\nu_{\mathrm{rel}} \to \eta/\rho_0 = \nu$, the ordinary
kinematic viscosity.

Similarly, the \emph{relativistic thermal diffusivity} is
\begin{equation}
\label{eq:rel2-25}
  \kappa_T = \frac{\kappa_{\mathrm{eff}}}{w_0\,c_p}\,,
\end{equation}
where $c_p$ is the specific heat at constant pressure.

%======================================================================
\section{The perturbation equations}\label{sec:rel9}
%======================================================================

%----------------------------------------------------------------------
\subsection{The equilibrium state}\label{sec:rel9a}
%----------------------------------------------------------------------

Consider an infinite horizontal layer of relativistic fluid in which a
steady adverse temperature gradient is maintained and no motions are present.
The equilibrium (unperturbed) state is characterized by
\begin{equation}
\label{eq:rel2-26}
  u^i = 0, \qquad T = T_0 - \beta z, \qquad Q^\mu = 0,
  \qquad \Pi^{\mu\nu} = 0, \qquad \Pi_{\mathrm{bulk}} = 0,
\end{equation}
where $\beta > 0$ is the adverse temperature gradient.  The equilibrium
pressure distribution satisfies the relativistic hydrostatic equation
\begin{equation}
\label{eq:rel2-27}
  \frac{\mathrm{d}p_0}{\mathrm{d}z} = -g\,w_0(z),
\end{equation}
with $w_0(z) = (\varepsilon_0 + p_0)/c^2$ evaluated at the local
temperature $T_0 - \beta z$.  In the Boussinesq approximation, the
enthalpy density varies linearly:
\begin{equation}
\label{eq:rel2-28}
  w_0(z) = w_0\bigl(1 + \alpha_w\,\beta\,z\bigr).
\end{equation}

%----------------------------------------------------------------------
\subsection{Perturbation of the energy--momentum tensor}\label{sec:rel9b}
%----------------------------------------------------------------------

Let the equilibrium state be slightly perturbed.  We write
\begin{equation}
\label{eq:rel2-29}
  u^i \to u^i, \quad T \to T_0 - \beta z + \theta, \quad
  p \to p_0(z) + \delta p,
\end{equation}
where $u^i$, $\theta$, and $\delta p$ are first-order perturbation
quantities.  In the BDNK formalism, the dissipative fluxes $Q^\mu$ and
$\Pi^{\mu\nu}$ are \emph{not} independent perturbation variables; they are
determined algebraically from the gradients of $(u^i, \theta, \delta p)$
via the constitutive relations~\eqref{eq:rel2-14}--\eqref{eq:rel2-17}.

The perturbation of the energy--momentum tensor~\eqref{eq:rel2-7}
at first order is
\begin{equation}
\label{eq:rel2-30}
  \delta T^{00} = \delta\varepsilon + \delta\varepsilon_1, \qquad
  \delta T^{0i} = w_0\,c\,u^i + \frac{\delta Q^i}{c}, \qquad
  \delta T^{ij} = \delta p\,\delta^{ij} + \delta\Pi^{ij},
\end{equation}
where $\delta\varepsilon = (\partial\varepsilon/\partial T)_p\,\theta
+ (\partial\varepsilon/\partial p)_T\,\delta p$, and
$\delta\varepsilon_1$ is the BDNK frame correction to the energy density
(first order in gradients of the perturbation).  The heat flux perturbation
$\delta Q^i$ and anisotropic stress perturbation $\delta\Pi^{ij}$ are
determined by the linearized constitutive relations.

%----------------------------------------------------------------------
\subsection{Linearized equations of motion}\label{sec:rel9c}
%----------------------------------------------------------------------

Linearizing the momentum equation~\eqref{eq:rel2-23} about the equilibrium
state and retaining only first-order terms, we obtain
\begin{equation}
\label{eq:rel2-31}
  \frac{\partial u_i}{\partial t}
  = -\frac{1}{w_0}\frac{\partial\,\delta p}{\partial x_i}
    + g\,\alpha_w\,\theta\,\lambda_i
    + \nu_{\mathrm{rel}}\,\nabla^2 u_i\,.
\end{equation}
This is the relativistic counterpart of equation~(55) in Chapter~II,
with $\nu_{\mathrm{rel}} = \eta/w_0$ replacing $\nu = \mu/\rho_0$ and
$\alpha_w$ replacing $\alpha$.

\medskip
\noindent\textbf{Remark.}\quad
In the Israel--Stewart formulation, the momentum equation at this stage
would couple to a separate relaxation equation for $\pi_{ij}$ involving
$\tau_\pi\,\partial\pi_{ij}/\partial t$.  In the BDNK formalism, the shear
stress $\Pi_{ij} = -2\eta\,\sigma_{ij}$ is an algebraic function of the
velocity gradients, so the substitution is immediate:
$\partial_j \Pi_{ij} = \eta\,\nabla^2 u_i$ (for incompressible flow).
The linearized momentum equation is therefore \emph{first order in time},
identical in structure to the classical Navier--Stokes equation.

%----------------------------------------------------------------------
\subsection{Linearized heat equation}\label{sec:rel9d}
%----------------------------------------------------------------------

Linearizing the energy equation~\eqref{eq:rel2-13} about the equilibrium
state and substituting the BDNK constitutive relation for the heat flux
from equation~\eqref{eq:rel2-14b}, we obtain the linearized heat equation
directly.

At first order, the divergence of $Q^\mu$ yields
\begin{equation}
  \nabla_\mu Q^\mu \big|_{\text{1st order}}
  = -\kappa_{\mathrm{eff}}\,\nabla^2\theta + \text{(BDNK frame corrections)}.
\end{equation}
Substituting into the energy balance gives
\begin{equation}
\label{eq:rel2-34}
  w_0\,c_p\frac{\partial\theta}{\partial t}
  = w_0\,c_p\,\beta\,w
    + \kappa_{\mathrm{eff}}\,\nabla^2\theta\,,
\end{equation}
where $w = \lambda_j u_j = u_3$ is the vertical component of velocity
and $c_p$ is the specific heat at constant pressure.  Dividing by
$w_0\,c_p$ and using $\kappa_T = \kappa_{\mathrm{eff}}/(w_0 c_p)$:
\begin{equation}
\label{eq:rel2-36}
  \boxed{
  \frac{\partial\theta}{\partial t}
  = \beta\,w + \kappa_T\,\nabla^2\theta\,.}
\end{equation}
This is the \textbf{BDNK heat equation} for the temperature perturbation.
It is the relativistic replacement for equation~(56) of Chapter~II and
is \emph{first order in time}---in contrast to the Israel--Stewart
telegraph equation
$\tau_q\,\partial^2\theta/\partial t^2 + \partial\theta/\partial t
= \beta(\tau_q\,\partial w/\partial t + w) + \kappa_T\nabla^2\theta$,
which is second order in time.

Equation~\eqref{eq:rel2-36} is identical in form to the classical
equation~(76) of Chapter~II; the relativistic content enters through the
definitions of $\kappa_T$ and the BDNK frame coefficients that ensure
causality.  The fact that BDNK yields a first-order-in-time heat equation
(rather than a telegraph equation) reflects the central feature of the
formalism: \emph{no relaxation times are introduced}.

%----------------------------------------------------------------------
\subsection{The solenoidal condition}\label{sec:rel9e}
%----------------------------------------------------------------------

From the relativistic Boussinesq approximation~\eqref{eq:rel2-19}
and the low-velocity limit, the velocity field remains solenoidal:
\begin{equation}
\label{eq:rel2-38}
  \frac{\partial u_j}{\partial x_j} = 0,
\end{equation}
as in the classical case (equation~(57) of Chapter~II).

%----------------------------------------------------------------------
\subsection{Perturbation equations for vorticity and vertical velocity}%
\label{sec:rel9f}
%----------------------------------------------------------------------

Following exactly the classical procedure of \S~9 (equations~(65)--(76) of
Chapter~II), we take the curl and double curl of the momentum
equation~\eqref{eq:rel2-31} to eliminate the pressure perturbation.

Defining the vorticity $\omega_i = \epsilon_{ijk}\,\partial u_k/\partial x_j$
and the vertical components $\zeta = \lambda_i\omega_i$,
$w = \lambda_i u_i$, we obtain:

\paragraph{Vorticity equation (vertical component):}
\begin{equation}
\label{eq:rel2-39}
  \frac{\partial\zeta}{\partial t} = \nu_{\mathrm{rel}}\,\nabla^2\zeta\,,
\end{equation}
which is identical in form to the classical equation~(73), but with
$\nu_{\mathrm{rel}} = \eta/w_0$ replacing $\nu = \mu/\rho_0$.

\paragraph{Vertical velocity equation:}
\begin{equation}
\label{eq:rel2-40}
  \frac{\partial}{\partial t}\nabla^2 w
  = g\,\alpha_w\!\left(
      \frac{\partial^2\theta}{\partial x_1^2}
      + \frac{\partial^2\theta}{\partial x_2^2}
    \right)
    + \nu_{\mathrm{rel}}\,\nabla^4 w\,,
\end{equation}
which is the relativistic counterpart of equation~(74), with
$g\alpha \to g\alpha_w$ (reflecting the enthalpy contribution to buoyancy)
and $\nu \to \nu_{\mathrm{rel}}$.

Together with the BDNK heat equation~\eqref{eq:rel2-36}, the
solenoidal condition~\eqref{eq:rel2-38}, and the vorticity
equation~\eqref{eq:rel2-39}, equations~\eqref{eq:rel2-40} constitute the
\emph{complete set of linearized perturbation equations for the
relativistic B\'enard problem in the BDNK formalism}.

\medskip
\noindent\textbf{Remark on the order of the system.}\quad
The BDNK perturbation system is \emph{lower order} than the Israel--Stewart
system: the heat equation~\eqref{eq:rel2-36} is first order in time
(not second), and the momentum equation~\eqref{eq:rel2-31} involves no
relaxation dynamics for $\Pi_{ij}$.  The total number of independent
dynamical equations is reduced because $Q^\mu$, $\Pi^{\mu\nu}$, and
$\Pi_{\mathrm{bulk}}$ are determined algebraically.  Despite being lower
order, the full BDNK system (before the Boussinesq/low-velocity limit)
remains \emph{strongly hyperbolic}, as established rigorously by
Bemfica, Disconzi, and Noronha.

%----------------------------------------------------------------------
\subsection{Boundary conditions}\label{sec:rel9g}
%----------------------------------------------------------------------

The boundary conditions on the perturbations are specified at $z = 0$ and
$z = d$.  For rigid, perfectly conducting boundaries:
\begin{equation}
\label{eq:rel2-41}
  w = 0, \qquad \frac{\partial w}{\partial z} = 0, \qquad
  \theta = 0 \qquad \text{at } z = 0, d.
\end{equation}
In the BDNK formalism, because the heat flux is determined algebraically
by the constitutive relation~\eqref{eq:rel2-14b}, no additional boundary
condition on $Q_z$ is needed beyond the temperature boundary condition
$\theta = 0$ (which already determines $Q_z$ through the gradient of
$\theta$).  This is in contrast to the Israel--Stewart formulation, where
the independent dynamical variable $q_z$ requires its own boundary
condition.

Alternatively, for ``free'' (stress-free) boundaries:
\begin{equation}
\label{eq:rel2-43}
  w = 0, \qquad \frac{\partial^2 w}{\partial z^2} = 0, \qquad
  \theta = 0 \qquad \text{at } z = 0, d.
\end{equation}
These are the same conditions as in the classical theory (cf.\ Chapter~II,
\S\S~12--14).

%----------------------------------------------------------------------
\subsection{Causality and hyperbolicity of the BDNK perturbation system}%
\label{sec:rel9h}
%----------------------------------------------------------------------

It is essential to verify that the relativistic perturbation equations
derived above respect causality, i.e.\ that perturbations propagate at
finite speed bounded by $c$.

\paragraph{Structure of the system.}
The complete linearized BDNK system consists of:
\begin{itemize}
\item Equation~\eqref{eq:rel2-31}: a \emph{first-order-in-time}
  equation for $u_i$, with the shear stress given algebraically by
  $\Pi_{ij} = -2\eta\,\sigma_{ij}$;
\item Equation~\eqref{eq:rel2-36}: a \emph{first-order-in-time}
  equation for $\theta$.
\end{itemize}
This is a \emph{lower-order} system than the Israel--Stewart perturbation
equations, which are second order in time due to the relaxation terms.

\paragraph{Causality via BDNK frame coefficients.}
In the BDNK formalism, causality is \emph{not} enforced by relaxation-time
bounds of the form $\tau_q \ge \kappa_T/c^2$ (as in Israel--Stewart).
Instead, causality is a consequence of the choice of hydrodynamic frame.
The system of conservation equations $\nabla_\mu T^{\mu\nu} = 0$,
$\nabla_\mu(\rho_0 u^\mu) = 0$, with the BDNK constitutive relations
\eqref{eq:rel2-7a}--\eqref{eq:rel2-7d}, is \textbf{strongly hyperbolic}
provided the transport coefficients and frame coefficients satisfy
the following coupled inequalities:

\begin{enumerate}
\item \textbf{Positive dissipation:}
\begin{equation}
\label{eq:rel2-44}
  \eta > 0, \qquad \zeta + \tfrac{2}{3}\eta > 0, \qquad \kappa > 0.
\end{equation}

\item \textbf{BDNK frame bounds:}
The frame coefficients ($\varepsilon_1$-type corrections, $\beta_1$,
$\alpha_1$, $\zeta_1$, etc.)\ must satisfy a set of coupled nonlinear
inequalities that ensure the principal symbol of the PDE system has real
eigenvalues and a complete set of eigenvectors.  In the notation of
Bemfica--Disconzi--Noronha--Kovtun (2023), these are:
\begin{equation}
\label{eq:rel2-45}
  \mathcal{A}(\eta, \zeta, \kappa, \beta_1, \varepsilon_1, \ldots) > 0,
  \qquad
  \mathcal{B}(\eta, \zeta, \kappa, \beta_1, \varepsilon_1, \ldots) > 0,
\end{equation}
where $\mathcal{A}$ and $\mathcal{B}$ denote specific combinations of the
transport and frame coefficients (see Bemfica et al., Phys.\ Rev.\ D 107,
076012, 2023, for the explicit forms).

\item \textbf{Subluminal signal speeds:} The characteristic speeds of
the system are bounded:
\begin{equation}
\label{eq:rel2-46}
  |v_{\mathrm{char}}| \le c,
\end{equation}
which is guaranteed when conditions~\eqref{eq:rel2-44}
and~\eqref{eq:rel2-45} are satisfied.
\end{enumerate}

\paragraph{Comparison with Israel--Stewart causality.}
In the Israel--Stewart theory, causality requires the relaxation-time bounds
$\tau_q \ge \kappa_T/c^2$ and $\tau_\pi \ge \nu_{\mathrm{rel}}/c^2$,
which ensure that the thermal and viscous signal speeds remain subluminal.
The BDNK approach replaces these conditions with the frame-coefficient
inequalities~\eqref{eq:rel2-45}, which are \emph{rigorous} (proven at
the fully nonlinear level, not merely in linearization) and do not
introduce any additional dynamical modes.

\paragraph{Dispersion relation.}
For plane-wave perturbations $\propto \exp[i(\mathbf{k}\cdot\mathbf{x} - \omega t)]$,
the BDNK dispersion relation is a \emph{lower-degree} polynomial in $\omega$
than the Israel--Stewart one (which has additional ``relaxation'' modes).
Specifically, the heat sector yields a standard diffusive dispersion
relation at low $k$,
\begin{equation}
\label{eq:rel2-47}
  \omega = -i\,\kappa_T\,k^2 + \mathcal{O}(k^3),
\end{equation}
while at high $k$ the BDNK frame coefficients modify the dispersion to
ensure $|\mathrm{Re}\,\omega/k| \le c$.  There are no spurious
high-frequency modes.

\paragraph{Conclusion.}
The BDNK perturbation system is \textbf{strongly hyperbolic}
(in the sense of Bemfica--Disconzi--Noronha), and all signal speeds are
bounded by $c$, provided the frame coefficients satisfy the coupled
inequalities~\eqref{eq:rel2-45}.  The classical parabolic limit is
recovered when $c \to \infty$ (in which case the frame corrections become
negligible and the system reduces to Fourier/Navier--Stokes theory).

\medskip
\noindent\textbf{Summary.}\quad
The BDNK perturbation equations for the B\'enard problem are:
\begin{alignat}{2}
\label{eq:rel2-summary}
  &\text{(i)}\quad
  &&\frac{\partial}{\partial t}\nabla^2 w
  = g\,\alpha_w\!\left(\frac{\partial^2}{\partial x_1^2}
    + \frac{\partial^2}{\partial x_2^2}\right)\theta
  + \nu_{\mathrm{rel}}\,\nabla^4 w, \\[4pt]
\label{eq:rel2-summary2}
  &\text{(ii)}\quad
  &&\frac{\partial\theta}{\partial t}
  = \beta\,w + \kappa_T\,\nabla^2\theta, \\[4pt]
\label{eq:rel2-summary3}
  &\text{(iii)}\quad
  &&\frac{\partial\zeta}{\partial t}
  = \nu_{\mathrm{rel}}\,\nabla^2\zeta, \\[4pt]
\label{eq:rel2-summary4}
  &\text{(iv)}\quad
  &&\frac{\partial u_j}{\partial x_j} = 0,
\end{alignat}
with $\nu_{\mathrm{rel}} = \eta/w_0$,
$\alpha_w = \alpha + (\partial p/\partial T)_\rho / (w_0 c^2)$,
$\kappa_T = \kappa_{\mathrm{eff}}/(w_0 c_p)$, and the BDNK causality
conditions on the frame coefficients
(equations~\eqref{eq:rel2-44}--\eqref{eq:rel2-45}).  In contrast to the
Israel--Stewart summary equations, the heat equation (ii) is first order in
time (no $\tau_q\,\partial^2\theta/\partial t^2$ term) and the shear stress
is algebraic (no $\tau_\pi$).

These reduce to the classical Chandrasekhar equations~(73), (74), (76), and
(57) of Chapter~II when $c \to \infty$ and $w_0 \to \rho_0$.
